\documentclass[12pt,a4paper]{article}

\usepackage[T2A]{fontenc}
\usepackage[utf8]{inputenc}
\usepackage[russian]{babel}
\usepackage[margin=22mm]{geometry}
\usepackage{enumitem}
\usepackage{array}

\setlength{\parindent}{0pt}
\setlength{\parskip}{5pt}
\setlist{nosep}

\begin{document}

\begin{titlepage}
\centering
\vspace*{20mm}

{\large НАЦИОНАЛЬНЫЙ ИССЛЕДОВАТЕЛЬСКИЙ УНИВЕРСИТЕТ\\
<<ВЫСШАЯ ШКОЛА ЭКОНОМИКИ>>\par}
\vspace{6mm}
{\large Факультет компьютерных наук\\
Департамент программной инженерии\par}

\vfill

{\Large\bfseries ОТЧЁТ\par}
\vspace{4mm}
{\Large по индивидуальному заданию №1\\
по дисциплине\\
<<Архитектура компьютера и операционные системы>>\par}

\vspace{12mm}
{\large\bfseries Вариант 29. Группа лифтов\par}

\vfill

\begin{flushright}
{\large Выполнил: Буров Иван Васильевич\\
Группа: БПИ252\par}
\end{flushright}

\vspace{12mm}
{\large Москва, 2026\par}
\end{titlepage}

%======================================================================
\section{Условие задачи}

В многоэтажном здании работает несколько лифтов под управлением общего
диспетчера. Пассажиры появляются на разных этажах через случайные
интервалы, каждый имеет номер и этаж назначения. Пассажир формирует вызов
с этажом отправления и направлением; диспетчер назначает вызов лифту по
выбранной стратегии. Лифт не принимает людей сверх вместимости: если
прибывшая кабина заполнена, оставшиеся продолжают ожидание, и вызов
считается выполненным только после формирования нового вызова для них.
Лифт обслуживает попутные вызовы; направление меняется после выполнения
заявок текущего направления.

\textbf{Входные параметры:} число этажей, лифтов и вместимость кабин;
начальное положение лифтов; количество пассажиров и их этажи; диапазон
времени появления; скорость движения и время дверей, посадки и высадки;
стратегия назначения вызовов.

\textbf{Инварианты:} вместимость не превышается; пассажир входит и выходит
только при открытых дверях на этаже остановки; лифт не движется с
открытыми дверями; пассажир находится либо на этаже, либо в одном лифте;
выходит только на заданном этаже; один вызов не назначается нескольким
лифтам.

\textbf{Завершение:} после доставки всех пассажиров либо по истечении
заданного периода; для оставшихся выводится состояние их заявок.

%======================================================================
\section{Отображение сущностей в программные конструкции}

Задание требует последовательной имитации, поэтому участники модели
отображаются не в процессы, а в структуры данных и функции, обрабатываемые
по очереди на каждом такте. Время дискретно; все длительности заданы в
тактах и являются входными параметрами.

\begin{center}
\begin{tabular}{@{}p{30mm}p{32mm}p{85mm}@{}}
\textbf{Сущность} & \textbf{Конструкция} & \textbf{Содержание} \\
\hline
Здание & \verb|Sim| & всё состояние: кабины, пассажиры, реестр вызовов,
время, счётчики \\
Лифт & \verb|Elevator| & этаж, направление, состояние автомата, таймер
фазы, признак открытых дверей, список едущих \\
Пассажир & \verb|Passenger| & номер, откуда, куда, направление, статус,
такты появления, посадки и доставки \\
Вызов & \verb|Call| & этаж, направление, назначенный лифт, такт
регистрации, признак активности \\
Диспетчер & \verb|dispatch| & собственного состояния не имеет: раздаёт
нераспределённые вызовы \\
\hline
\end{tabular}
\end{center}

\textbf{Один такт} (функция \verb|sim_step|): появление очередного
пассажира $\to$ назначение нераспределённых вызовов $\to$ обновление
каждой кабины $\to$ проверка инвариантов $\to$ отрисовка и пауза.

\textbf{Жизненный цикл вызова.} Пассажир появляется и либо регистрирует
вызов, либо присоединяется к существующему: на этаж и направление активен
не более одного вызова. Диспетчер назначает его одной кабине и больше не
переназначает. Кабина приезжает и забирает \emph{всех} ожидающих,
едущих в её сторону, пока есть место, в том числе тех, чей вызов был
назначен другой кабине --- это и есть обслуживание попутных вызовов. Вызов
гасится, а для не поместившихся немедленно регистрируется новый.

%======================================================================
\section{Структура программы и автомат кабины}

Программа написана на C, собирается командой \verb|make| с ключами
\verb|-std=c17 -D_POSIX_C_SOURCE=200809L -Wall -Wextra -O2|, предупреждений
нет. Макрос \verb|_POSIX_C_SOURCE| задан явно: при строгом \verb|-std=c17|
glibc не объявила бы \verb|dprintf|.

\begin{center}
\begin{tabular}{@{}p{28mm}r p{95mm}@{}}
\textbf{Файл} & \textbf{Строк} & \textbf{Содержимое} \\
\hline
\verb|elevator.h| & 102 & типы, константы, объявления \\
\verb|config.c| & 250 & умолчания, конфигурационный файл, аргументы,
проверка значений \\
\verb|sim.c| & 396 & пассажиры, вызовы, диспетчер, автомат кабины,
инварианты \\
\verb|output.c| & 180 & схема здания, лента событий, журнал, статистика \\
\verb|main.c| & 89 & сигналы, главный цикл, причина завершения \\
\hline
\end{tabular}
\end{center}

Весь ввод-вывод идёт через файловые дескрипторы: \verb|open|,
\verb|read|, \verb|write|, \verb|close|, \verb|dprintf|. Потоки
\verb|FILE*| не используются нигде. Также применяются \verb|sigaction|
и \verb|nanosleep|.

\textbf{Автомат кабины} состоит из шести состояний: \verb|ST_IDLE|,
\verb|ST_MOVING| (длительность \verb|travel_ticks|), \verb|ST_OPENING| и
\verb|ST_CLOSING| (\verb|door_ticks|), \verb|ST_UNLOADING| и
\verb|ST_LOADING| (длительность пропорциональна числу входящих или
выходящих). Такое устройство само по себе обеспечивает три инварианта:
пассажир не может войти или выйти вне состояний посадки и высадки, а
кабина не может ехать, не находясь в \verb|ST_MOVING|.

\textbf{Выбор направления.} Функция \verb|work_in_direction| отвечает,
есть ли ещё причина ехать в текущую сторону: пассажир внутри, которому
нужен этаж в этом направлении, либо назначенный кабине вызов с такого
этажа. Пока причина есть, направление не меняется --- это требование
условия. Отдельного массива целевых этажей у кабины нет: вся информация
уже содержится в списке едущих и реестре вызовов, а дублирующий массив
пришлось бы синхронно поддерживать.

%======================================================================
\section{Диспетчер и проверка инвариантов}

Диспетчер на каждом такте просматривает реестр и выбирает кабину для
каждого нераспределённого вызова. Заполненные кабины в выборе не
участвуют; если свободных нет, вызов рассматривается на следующем такте.

\textbf{\texttt{fcfs}} --- вызовы обслуживаются в порядке поступления:
берётся первая кабина без пассажиров и назначенных вызовов, а если все
заняты --- наименее загруженная.

\textbf{\texttt{nearest}} --- берётся кабина, которая быстрее доберётся до
этажа вызова. Оценка --- расстояние в этажах, кабине, едущей в
противоположную сторону, добавляется штраф в высоту здания.

Назначенный вызов никогда не переназначается. Это осознанное решение:
именно оно делает невозможным нарушение шестого инварианта. Платой
является неоптимальность --- освободившаяся рядом кабина не перехватит
вызов, назначенный дальней.

\textbf{Инварианты} проверяются во время работы; нарушение выводится
строкой \verb|INVARIANT BROKEN| и даёт код возврата 2. Четыре из них
проверяются после каждого такта в \verb|check_invariants|: вместимость,
закрытые двери при движении, нахождение пассажира ровно в одной кабине,
отсутствие двух активных вызовов на этаж и направление. Два оставшихся ---
вход и выход только при открытых дверях и выход только на своём этаже ---
проверяются в \verb|do_boarding| и \verb|do_unloading|, поскольку это
единственные места, где пассажир меняет местоположение; по состоянию в
конце такта проверить их уже невозможно.

%======================================================================
\section{Пользовательский интерфейс}

Экран перерисовывается после каждого такта: сверху схема здания, снизу
лента десяти последних событий.

{\small
\begin{verbatim}
tick 45    delivered 3/25   strategy nearest   in the building 4

  fl   E1     E2     E3    hall
  12    .      .      .
  11    .      .      .
  10    .      .      .   1v
   9    .      .    [^0]
   8    .      .      .
   7    .      .      .
   6    .      .      .
   5    .      .      .
   4  [v1]     .      .   1^
   3    .      .      .
   2    .    [^1]     .
   1    .      .      .

[  43] passenger 7 appears on floor 4 and wants floor 5 (up)
[  43] call registered on floor 4, up
[  43] dispatcher sends elevator 1 to the call on floor 4 (up)
[  44] elevator 1 starts going down from floor 4
[  45] elevator 3 reaches floor 9
\end{verbatim}
}

Обозначения: \verb|[^2]| --- кабина с закрытыми дверями едет вверх,
внутри двое; \verb|[v2]| --- вниз, \verb|[-2]| --- стоит;
\verb|(^2)| --- двери открыты; точка --- шахта на этаже пуста; столбец
\verb|hall| показывает, сколько человек ждёт и куда. Кабина в движении
выделяется жёлтым, с открытыми дверями --- зелёным.
Скорость задаётся параметром \verb|--delay-ms|, по умолчанию 200 мс. Флаг
\verb|--plain| отключает перерисовку и печатает события сплошной лентой:
этот режим нужен для тестовых скриптов и сравнения журналов.

Журнал пишется параллельно экрану через отдельный дескриптор и содержит
шапку с параметрами, все события и итоговую статистику.

%======================================================================
\section{Варианты запуска и завершения}

Параметры берутся из трёх источников с возрастающим приоритетом:
встроенные умолчания, конфигурационный файл, аргументы командной строки.
Оба источника обрабатываются одной функцией \verb|set_option|, поэтому
набор ключей у них одинаков; в файле они пишутся через подчёркивание, в
командной строке через дефис. Полный список выводится по \verb|--help|.
Все значения проверяются до начала моделирования: при недопустимом
выводится диапазон и программа завершается с кодом 1, не начав работу.

{\small
\begin{verbatim}
./build/elevator                                      # умолчания
./build/elevator --config data/rush-hour.conf         # сценарий из файла
./build/elevator --config data/default.conf --floors 20 --strategy fcfs
./build/elevator --plain --delay-ms 0 --seed 42 --log run.log
\end{verbatim}
}

Предусмотрены три способа завершения, и все приводят к одинаковому финалу:
выводится статистика и состояние незакрытых заявок.

\begin{enumerate}
\item \textbf{Все пассажиры доставлены} --- естественное завершение.
\item \textbf{Исчерпан лимит тактов} \verb|--max-ticks| --- когда
      интересует поведение системы за фиксированный период.
\item \textbf{Прерывание пользователем} --- SIGINT и SIGTERM
      перехватываются через \verb|sigaction|. Обработчик выставляет
      единственную переменную \verb|volatile sig_atomic_t| и больше ничего
      не делает. Главный цикл видит флаг, выходит, восстанавливает курсор
      терминала, печатает статистику и закрывает журнал.
\end{enumerate}

Коды возврата: 0 --- норма, 1 --- ошибка в параметрах, 2 --- нарушен
инвариант модели.

{\small
\begin{verbatim}
--- simulation finished: interrupted by the user ---
model time            41 ticks
delivered             2 of 25
average waiting       8.5 ticks
longest waiting       12 ticks
invariants broken     0
elevator 1            12 floors passed, 0 on board at the end

requests still open:
  passenger 3 waits on floor 7 for floor 3 since tick 14
  passenger 4 rides elevator 2 towards floor 6
  17 passenger(s) never appeared
\end{verbatim}
}

%======================================================================
\section{Тестирование и результаты}

В каталоге \verb|data/| подготовлено пять сценариев, каждый является
одновременно примером конфигурации и тестовым набором: \verb|small| ---
5 этажей и один лифт; \verb|default| --- 12 этажей, 3 лифта, вместимость
6; \verb|rush-hour| --- час пик, 16 этажей, 4 лифта, пассажир каждые
1--2 такта; \verb|slow-doors| --- то же здание, что базовое, но втрое
медленнее двери и посадка; \verb|crowded| --- вместимость 1, почти каждый
вызов приходится регистрировать заново.

Скрипт \verb|./tests/run-tests.sh| пересобирает проект и выполняет 30
проверок: сборка без предупреждений; отказ на семи вариантах неверного
ввода; прохождение всех пяти сценариев; проверка, что \verb|crowded|
действительно доводит до переполнения кабин; обе стратегии; краевые случаи
(ноль пассажиров, два этажа при восьми лифтах, один лифт вместимостью
один, обрыв по лимиту тактов); воспроизводимость по зерну в обе стороны;
совпадение журнала с экраном; обработка SIGINT; 60 прогонов со случайными
конфигурациями здания. Скрипт возвращает 0, только если прошли все.

Генератор случайных чисел реализован в программе собственный, а не взят из
библиотеки: \verb|rand| не обязана давать одинаковую последовательность на
разных системах, а для сравнения журналов между прогонами это необходимо.

\begin{center}
\begin{tabular}{@{}lrrrrrr@{}}
\textbf{Сценарий} & \textbf{Тактов} & \textbf{Доставлено} &
\textbf{Ожидание} & \textbf{Макс.} & \textbf{Поездка} & \textbf{Отказов} \\
\hline
\verb|small| & 65 & 8/8 & 15,0 & 27 & 8,6 & 0 \\
\verb|default| & 191 & 25/25 & 13,4 & 36 & 14,8 & 0 \\
\verb|rush-hour| & 151 & 60/60 & 17,2 & 59 & 16,8 & 0 \\
\verb|slow-doors| & 227 & 25/25 & 43,1 & 100 & 31,4 & 0 \\
\verb|crowded| & 223 & 20/20 & 70,2 & 206 & 12,8 & 21 \\
\hline
\end{tabular}
\end{center}

Времена указаны в тактах. Результаты согласуются со здравым смыслом: в
\verb|slow-doors| здание и поток те же, что в базовом сценарии, но все
длительности увеличены втрое --- среднее ожидание выросло с 13,4 до 43,1
такта. В \verb|crowded| вместимость равна одному человеку, ожидание
подскакивает до 70 тактов при 21 отказе в посадке, тогда как сама поездка
остаётся самой короткой: здание низкое.

Сравнение стратегий на одном потоке пассажиров (\verb|--seed 42|, базовое
здание): \verb|fcfs| --- 145 тактов, среднее ожидание 14,3;
\verb|nearest| --- 135 тактов, среднее ожидание 15,2. Стратегия
\verb|nearest| завершает обслуживание раньше, но среднее ожидание у неё
чуть выше: отправляя к вызову ближайшую кабину, она группирует кабины в
одной части здания, и дальние вызовы ждут дольше. На одном прогоне делать
выводы некорректно, для содержательного сравнения нужно усреднение по
многим зёрнам.

Во всех 60 прогонах со случайными конфигурациями (2--26 этажей, 1--8
лифтов, вместимость 1--6) все пассажиры доставлены, счётчик нарушений
инвариантов остался нулевым.

%======================================================================
\section{Выводы}

Разработано консольное приложение, моделирующее работу группы лифтов под
управлением общего диспетчера: один процесс, один поток, дискретное
модельное время. Реализованы все требования условия --- вызовы с этажом и
направлением, две стратегии диспетчеризации, ограничение вместимости с
перерегистрацией вызова для не поместившихся, обслуживание попутных
вызовов, смена направления только после отработки заявок текущего
направления. Все шесть инвариантов проверяются во время работы.
Предусмотрены три способа завершения, включая прерывание пользователем, и
во всех случаях выводится состояние незакрытых заявок. Корректность
подтверждена набором из 30 автоматических проверок.

\section{Использовани ИИ}

Использовал нейросеть Claude (Anthropic) для помощи в организации архитектуры проекта, добавления комментариев, составления тестов и отчета.

\end{document}
